\documentclass[12pt]{article}
\usepackage{amsmath, amssymb, amsthm, bm}
\usepackage[margin=1in]{geometry}
\usepackage[skip=0.6em]{parskip}

\newtheorem{proposition}{Proposition}
\theoremstyle{remark}
\newtheorem{remark}{Remark}

\begin{document}

\begin{center}
{\Large Moment selection and the Gaussian limit experiment}\\[3pt]
{\small Ivan A. Canay :: August 28, 2026}
\end{center}

\medskip

This note clarifies the relation between finite-sample generalized moment
selection and the Gaussian limit experiment used to study local power and
admissibility.  The main conclusion is simple.  The fixed-drift Gaussian
experiment is the correct limit experiment for the question of interest.
Standard moment selection retains every inequality whose standardized drift
has a finite limit.  It deletes a moment only when its slackness diverges on
the slower selection scale.  Thus moment selection matters at the boundary of
the local experiment, not at an ordinary finite point of that experiment.

This distinction also clarifies the multiple-core problem.  Deleting an
inequality that is far inside the null is different from determining which
core moment realizes the maximum in the local mean restriction.  The latter
cannot be learned without error when the core means have a finite local gap.

\section{A convenient normalization}

Suppose that the null moment inequalities are
\[
E[m_j(W,\theta)]\le 0,
\qquad j=1,\ldots,d.
\]
Define the studentized statistics and their population drifts by
\[
X_{j,n}
:=
\frac{\sqrt n\,\overline m_{j,n}(\theta_n)}
     {\widehat\sigma_{j,n}},
\qquad
\mu_{j,n}
:=
\frac{\sqrt n\,E[m_j(W,\theta_n)]}
     {\sigma_{j,n}}.
\]
Positive values indicate violation and negative values indicate slackness.
Under the usual local asymptotic approximation,
\begin{equation}\label{eq:gaussian-limit}
X_n\ \Rightarrow\ X\sim N_d(\mu,\Omega)
\qquad\text{when}\qquad
\mu_n\to\mu\in\mathbf R^d,
\end{equation}
where $\Omega$ is the limiting correlation matrix.

For concreteness, consider the hard-threshold version of generalized moment
selection.  Let $\kappa_n$ satisfy
\begin{equation}\label{eq:kappa}
\kappa_n\to\infty,
\qquad
\frac{\kappa_n}{\sqrt n}\to0,
\end{equation}
and retain moment $j$ when
\[
X_{j,n}\ge-\kappa_n.
\]
Equivalently, the selected set is
\begin{equation}\label{eq:selected-set}
\widehat{\mathcal J}_n
:=
\{j:X_{j,n}\ge-\kappa_n\}.
\end{equation}
Other GMS functions lead to the same conclusion below because their argument
is $X_n/\kappa_n$, which converges to zero under every finite local drift.

\section{Two asymptotic scales}

The root-$n$ drift $\mu_{j,n}$ governs the limiting statistic and local power.
The ratio $\mu_{j,n}/\kappa_n$ governs moment selection.  These are different
asymptotic objects.

\begin{proposition}[What GMS selects]\label{prop:gms-scales}
Suppose that
\[
X_n=\mu_n+Z_n+o_p(1),
\qquad
Z_n=O_p(1),
\]
and let $\widehat{\mathcal J}_n$ be defined by
\eqref{eq:selected-set}.

\begin{enumerate}
\item If $\mu_n\to\mu\in\mathbf R^d$, then
\[
P(\widehat{\mathcal J}_n=\{1,\ldots,d\})\to1.
\]

\item If, for some $j$,
\[
\frac{\mu_{j,n}}{\kappa_n}\to-c
\]
with $c>1$, then moment $j$ is deleted with probability tending to one.  If
$c<1$, it is retained with probability tending to one.

\item If the $j$th population inequality is separated from equality by a
fixed amount, so that
\[
E[m_j(W,\theta_n)]\le-\delta
\]
for some $\delta>0$ and its standard deviation remains bounded, then moment
$j$ is deleted with probability tending to one.
\end{enumerate}
\end{proposition}

\begin{proof}
Since $Z_n=O_p(1)$ and $\kappa_n\to\infty$,
\[
\frac{X_n}{\kappa_n}
=
\frac{\mu_n}{\kappa_n}+o_p(1).
\]
If $\mu_n$ has a finite limit, every coordinate of this ratio converges to
zero, which proves the first claim.  Under the second claim,
$X_{j,n}/\kappa_n\to_p-c$, so the selection rule is eventually on the
appropriate side of $-1$.  For the third claim,
$\mu_{j,n}$ is negative and of order $\sqrt n$.  Condition~\eqref{eq:kappa}
then gives $\mu_{j,n}/\kappa_n\to-\infty$.
\end{proof}

The boundary case $c=1$ is deliberately excluded from the second part of the
proposition.  At that boundary, selection depends on the centering at a finer
scale and need not have a unique limit.  This is the role of the additional
selection-scale parameter in \texttt{Addition-2.pdf}.

Proposition~\ref{prop:gms-scales} distinguishes three uses of the word
``slack.''  A moment with a finite negative value of $\mu_j$ is locally slack,
but it is statistically close to a binding moment and GMS retains it.  A
moment with drift of order $-\kappa_n$ lies on the selection scale.  A moment
that is slack by a fixed population amount has drift of order $-\sqrt n$ and
is deleted asymptotically.

\section{Why the Gaussian experiment does not contain an oracle selector}

Moment selection is a component of a testing rule, not a primitive of a
statistical experiment.  One can define a data-dependent selector in the
Gaussian experiment~\eqref{eq:gaussian-limit}, but it cannot identify the
binding set without error over the finite mean space.

\begin{proposition}[No oracle selection at finite drifts]
Let $X\sim N_d(\mu,\Omega)$ with $\Omega$ positive definite.  Suppose that two
finite mean vectors $\mu$ and $\nu$ have different binding-set or active-branch
labels.  There is no measurable selector that returns the correct label with
probability one under both $N_d(\mu,\Omega)$ and $N_d(\nu,\Omega)$.
\end{proposition}

\begin{proof}
Gaussian distributions with a common positive-definite covariance matrix and
finite mean vectors are mutually absolutely continuous.  Therefore, an event
that has probability one under $N_d(\mu,\Omega)$ also has probability one
under $N_d(\nu,\Omega)$.  If the two correct labels differ, the two events on
which a selector returns the correct label are disjoint.  They cannot both
have probability one.
\end{proof}

This is the limit-experiment version of the contiguity obstruction.  Local
slackness cannot be consistently distinguished from binding, and two active
branches separated by a finite local gap cannot be identified perfectly.
The limiting version of standard GMS avoids this classification problem by
retaining all finite-drift moments.

\section{The two-core problem}

Consider the identity-covariance experiment
\[
X\sim N_3(\mu,I_3),
\qquad
\mu_3\le\max\{\mu_1,\mu_2\},
\]
where coordinates 1 and 2 are core moments and coordinate 3 is redundant.
There are two distinct selection questions.

First, GMS asks whether an inequality is sufficiently far inside the null to
be deleted.  By Proposition~\ref{prop:gms-scales}, both core moments and the
redundant moment are retained whenever their local means are finite.

Second, a branchwise gated repair would need to determine which core realizes
$\max\{\mu_1,\mu_2\}$.  This is not the same question.  For example, the
natural rule that selects core 1 when $X_1>X_2$ satisfies
\begin{equation}\label{eq:active-core-probability}
P_\mu(X_1>X_2)
=
\Phi\left(\frac{\mu_1-\mu_2}{\sqrt2}\right).
\end{equation}
If core 1 is active, the expression in~\eqref{eq:active-core-probability} is
the probability of selecting it.  This probability equals $1/2$ at a tie and
is strictly below one for every finite core gap.  It converges to one only
when the standardized gap diverges.

Moreover, a core moment can fail to realize the maximum while still having a
positive local mean.  Such a moment is locally violated, not slack relative
to the null, and GMS will not delete it.  Thus ordinary GMS cannot be
interpreted as an asymptotically perfect active-core selector.

When one core is slack on the GMS scale, the situation is different.  That
core is deleted with probability tending to one, and the remaining problem
reduces to the one-core experiment.  The bivariate gated improvement can then
apply to that separated subproblem.  This is a lower-dimensional boundary
regime, not a solution to the problem over the full finite-drift mean space.

\section{Implications for admissibility}

Admissibility is defined for a fixed experiment and a fixed parameter space.
For local asymptotic admissibility, the relevant experiment is
\[
X\sim N_d(\mu,\Omega),
\qquad
\mu\in\Lambda\subseteq\mathbf R^d,
\]
with the full vector of finite local means retained.  There are three reasons
this is the natural object.

\begin{enumerate}
\item It is the weak limit of the vector of studentized sample moments along
root-$n$ local sequences.

\item It is the fixed-drift limit of the GMS procedure itself: the selection
input converges to zero, so the limiting critical value treats every moment
as potentially binding.

\item The alternatives for which moment selection deletes a coordinate lie
on a different scale.  They correspond to sequences escaping to the boundary
of the finite mean space rather than to omitted points of the ordinary local
experiment.
\end{enumerate}

Consequently, the sensitivity of the long or max test to a finite negative
mean is not an artifact.  It is exactly the local behavior that the limit
experiment is intended to record.  Moment selection recovers the appropriate
lower-dimensional problem far from the local boundary, but it does not remove
the multiple-core adaptation problem within that boundary.

In particular, the unresolved admissibility question in
\texttt{three-dimensional-case.pdf} remains genuine.  A branch-specific
bivariate dominator is not a single test that dominates over the union of the
two core branches.  Any global repair must account for the nonzero probability
of selecting the wrong active core when the core gap is finite.

\section{A useful scale-separation statement}

The connection between the finite-sample and limit analyses can be summarized
as follows.

\begin{enumerate}
\item On every fixed finite-drift sequence, and uniformly on compact drift
sets under the corresponding uniform central limit theorem, GMS retains all
moments.  The full-dimensional Gaussian test is therefore the relevant object
for local power and admissibility.

\item On sequences where an inequality is slack by more than the GMS
threshold, that moment is deleted and the limiting testing problem becomes
lower-dimensional.

\item Near a multiple-core branch boundary, no selector can determine the
active core without nontrivial error.  This transition region is where the
three-dimensional admissibility problem has bite.
\end{enumerate}

The two-rate analysis in \texttt{Addition-2.pdf} supplements the Gaussian
limit experiment by describing how the procedure behaves as parameters escape
the finite local mean space.  It does not replace that experiment.  For the
admissibility project, the fixed-drift Gaussian experiment is the correct and
natural starting point.

\end{document}
